\chapter{Introduction}
\label{ch:introduction}

\section{Motivation and Context}
\label{sec:motivation}

Financial markets are dynamic systems whose statistical properties change
substantially over time. Returns that are approximately stationary during
extended periods of economic expansion can suddenly give way to regimes
characterized by negative drift, elevated volatility, and strong
cross-sectional co-movement across assets and sectors. The global
financial crisis of 2008 and the COVID-19 market crash of March 2020 are
two of the most salient examples of this phenomenon in recent history. In
both episodes, equity indices such as the S\&P 500 experienced drawdowns
exceeding 30\% within a few weeks, volatility indices reached historical
highs, and the correlation structure between sectors that are normally
weakly related --- for instance, \textit{Energy} and \textit{Information
Technology} --- increased sharply, a hallmark of what the literature
broadly refers to as a \emph{crisis regime} \cite{ang2002regime}.

The notion that financial markets alternate between a small number of
distinct \emph{regimes} --- commonly described informally as
\emph{bull}, \emph{bear} and \emph{crisis} markets --- has a long
tradition in econometrics and quantitative finance. The seminal work of
Hamilton \cite{hamilton1989} introduced Markov-switching autoregressive
models to formally characterize shifts in the mean and variance of
macroeconomic and financial time series, providing the theoretical
foundation for a large body of subsequent research on
\emph{regime-switching models} \cite{guidolin2011}. From a more applied
perspective, asset managers, risk officers and algorithmic trading
systems routinely adapt their strategies depending on the prevailing
regime: position sizing, hedging ratios, and even the choice of which
factors to emphasize in a portfolio are often regime-dependent. Detecting
regime changes \emph{early} and \emph{reliably} is therefore directly
linked to the ability to manage risk and preserve capital during adverse
periods.

A second, more recent strand of research has explored the role of
\emph{sentiment} --- the aggregate tone of news, social media, or analyst
reports --- as a driver of, or at least a correlate of, market dynamics.
Early work in this area relied on dictionary-based methods, most notably
the finance-specific word lists developed by Loughran and McDonald
\cite{loughranmcdonald2011}, which showed that general-purpose sentiment
dictionaries (e.g., the Harvard-IV psychosocial dictionary) systematically
misclassify financial text, and that domain-specific lexicons are
necessary to obtain meaningful sentiment signals from corporate filings
and financial news. Other influential work, such as Tetlock's analysis of
the \emph{Wall Street Journal}'s ``Abreast of the Market'' column
\cite{tetlock2007} and Baker and Wurgler's investor sentiment index
\cite{baker2006}, established that aggregate measures of media tone carry
predictive information about subsequent market returns and volatility,
even after controlling for standard risk factors.

The advent of \emph{Transformer}-based language models
\cite{vaswani2017attention} and, in particular, of pre-trained
bidirectional encoders such as BERT \cite{devlin2019bert}, has
dramatically changed the toolkit available for sentiment extraction.
Fine-tuned, domain-adapted variants such as \textbf{FinBERT}
\cite{finbert2019} have been shown to substantially outperform both
lexicon-based methods and general-purpose sentiment classifiers on
financial text, producing well-calibrated probability estimates over
\emph{positive}, \emph{negative} and \emph{neutral} sentiment classes for
short pieces of financial text such as news headlines. At the same time,
the public release of large-scale, time-stamped financial news corpora
such as \textbf{FNSPID} (Financial News and Stock Price Integration
Dataset) \cite{zhang2023fnspid} has made it practical, for the first time
at this scale, to systematically compute sentiment scores for millions of
news items spanning more than a decade and to align them with daily stock
price data.

This confluence of factors --- a long-standing econometric tradition of
regime-switching models, a growing body of evidence that sentiment
contains information about market dynamics, and the availability of
both large financial news corpora and pre-trained finance-specific
language models --- motivates the central question addressed by this
thesis: \emph{can sector-level sentiment extracted from financial news
using FinBERT improve the identification of market regimes obtained
through unsupervised clustering of technical indicators, and if so, in
what sense?}

\section{Problem Statement}
\label{sec:problem-statement}

Most existing approaches to market regime detection rely either
exclusively on price-derived technical variables (returns, realized
volatility, momentum, moving-average ratios) processed with
regime-switching models or clustering algorithms such as \textit{K-Means}
\cite{ang2002regime,guidolin2011}, or exclusively on sentiment indices
constructed from news or social media, used as standalone predictors of
returns or volatility \cite{tetlock2007,baker2006}. Comparatively little
work has systematically combined both information sources within a
single \emph{unsupervised} clustering framework and then evaluated the
\emph{economic} content of the resulting groups --- as opposed to purely
statistical measures of cluster quality such as the silhouette
coefficient.

This thesis addresses this gap by:

\begin{enumerate}
  \item Building, from raw data, a reproducible pipeline that produces a
  daily dataset combining S\&P 500 technical indicators with sector-level
  (GICS) sentiment scores derived from FNSPID news headlines via FinBERT.
  \item Training and comparing three \textit{K-Means} clustering
  configurations that use, respectively, only technical indicators, both
  technical indicators and raw sector sentiment, and only smoothed sector
  sentiment.
  \item Training a Gaussian Hidden Markov Model on the same feature sets
  as a methodologically distinct comparison point, given that HMMs
  explicitly model temporal transition dynamics between regimes, unlike
  \textit{K-Means}.
  \item Evaluating all resulting partitions not only through internal
  clustering metrics (silhouette coefficient in a common
  principal-component space, feature importance via ANOVA F-statistics),
  but also through an \emph{economic validation} layer that computes
  standard portfolio performance metrics --- annualized return,
  volatility, Sharpe ratio, maximum drawdown, percentage of positive days,
  and extreme daily returns --- for every combination of GICS sector and
  detected regime.
\end{enumerate}

\section{Objectives}
\label{sec:objectives}

\subsection{General objective}

The general objective of this thesis is to \textbf{develop and validate a
clustering-based methodology for market regime detection, assessing
whether the incorporation of sector-level sentiment extracted from
financial news via FinBERT improves the economic differentiation of the
detected regimes relative to a purely technical baseline}.

\subsection{Specific objectives}

\begin{enumerate}
  \item \textbf{Data engineering.} Construct a unified, leakage-free
  dataset that merges (i) daily technical indicators computed from S\&P
  500 OHLC prices, and (ii) sector-level sentiment scores derived from the
  FNSPID news corpus, restricted to companies that were S\&P 500
  constituents at the time of publication and aligned to the following
  open market day.
  \item \textbf{Sentiment extraction.} Apply the FinBERT model to every
  filtered news headline to obtain a continuous sentiment score and a
  categorical sentiment label, and aggregate these scores at the
  (date, GICS sector) level using both raw daily statistics (mean and
  standard deviation) and a 21-session moving-average smoothing scheme.
  \item \textbf{Clustering model design.} Train three \textit{K-Means}
  models --- technical-only (Model A), technical + raw sentiment (Model
  B), and smoothed-sentiment-only (Model C) --- selecting the number of
  clusters $K=3$ via the elbow method and the silhouette coefficient, and
  label the resulting clusters as \emph{Bull}, \emph{Bear} or
  \emph{Crisis} based on centroid analysis.
  \item \textbf{Secondary comparison method.} Train a Gaussian Hidden
  Markov Model with $K=3$ states and full covariance matrices on the same
  feature sets used for Models A and B, and compare its temporal and
  economic properties with those of the \textit{K-Means} models.
  \item \textbf{Internal validation.} Compare the three \textit{K-Means}
  models using the silhouette coefficient computed in a common
  principal-component space (retaining 80\% of explained variance), 2D
  PCA visualizations, correlation analysis of the input features, and
  ANOVA-based feature importance.
  \item \textbf{Economic validation.} For every GICS sector and detected
  regime, compute annualized return, annualized volatility, Sharpe ratio
  ($r_f = 2\%$), maximum drawdown, percentage of positive days, and the
  best and worst single-day returns, and analyze the resulting heatmaps to
  characterize each regime economically.
  \item \textbf{Critical comparison.} Synthesize the results of Models
  A, B and C, and of the HMM comparison, into a coherent discussion of the
  incremental value (or lack thereof) of sentiment information for market
  regime detection.
\end{enumerate}

\section{Alignment with the Sustainable Development Goals}
\label{sec:sdgs}

In line with the commitment of Universidad Pontificia Comillas to the
2030 Agenda for Sustainable Development \cite{un2015sdgs}, this project
can be related to the following Sustainable Development Goals (SDGs):

\begin{itemize}
  \item \textbf{SDG 8 -- Decent Work and Economic Growth.} Tools that
  improve the understanding of market risk and instability contribute,
  indirectly but meaningfully, to more resilient financial systems. Early
  and more interpretable identification of \emph{crisis} regimes can help
  institutional investors and asset managers reduce unnecessary exposure
  during periods of systemic stress, contributing to more stable economic
  growth.
  \item \textbf{SDG 9 -- Industry, Innovation and Infrastructure.} The
  project applies state-of-the-art natural language processing
  (Transformer-based language models) and machine learning (unsupervised
  clustering, Hidden Markov Models) to a real-world financial data
  pipeline, illustrating how modern AI infrastructure can be integrated
  into quantitative finance workflows in a fully reproducible and
  open-source manner.
  \item \textbf{SDG 17 -- Partnerships for the Goals.} The entire pipeline
  is built on publicly available data (FNSPID, S\&P 500 historical
  constituents and prices) and open-source software (Python, pandas,
  scikit-learn, Hugging Face Transformers, hmmlearn), illustrating how
  open data and open-source tooling --- often the product of broad,
  decentralized collaboration --- can be leveraged to produce rigorous
  quantitative research without requiring access to proprietary data
  vendors.
\end{itemize}

\section{Structure of the Document}
\label{sec:structure}

The remainder of this document is organized as follows.

\textbf{\Cref{ch:soa} (State of the Art)} reviews the relevant literature
in two areas: market regime detection, covering Markov-switching models,
Hidden Markov Models and clustering-based approaches such as
\textit{K-Means}; and sentiment analysis in finance, covering
dictionary-based methods, general-purpose deep learning sentiment models,
and finance-specific models such as FinBERT, together with the FNSPID
dataset. The chapter closes with an identification of the gap addressed
by this thesis.

\textbf{\Cref{ch:methodology} (Methodology)} describes, in detail, the
full pipeline implemented for this project: the data sources and
preprocessing steps, the FinBERT-based sentiment scoring procedure, the
sector-level aggregation and smoothing scheme, the technical indicators
computed from S\&P 500 prices, the design of Models A, B and C
(\textit{K-Means}), the Gaussian HMM used as a secondary comparison
method, and the procedure used to label and interpret the resulting
clusters.

\textbf{\Cref{ch:results} (Results)} presents and discusses the results
obtained: internal clustering quality (silhouette coefficients, PCA
visualizations), the temporal distribution of the detected regimes,
correlation analysis of the input features, feature importance, the
economic validation heatmaps for Models A, B and C, and a comparison
between \textit{K-Means} and the Gaussian HMM.

\textbf{\Cref{ch:conclusions} (Conclusions and Future Work)} synthesizes
the main findings of the thesis, discusses the limitations of the proposed
methodology, and outlines several directions for future work.

Finally, the document includes a \textbf{Bibliography} with the
references cited throughout the text, and two \textbf{Annexes}: Annex A
describes the structure of the source code implementing the pipeline, and
Annex B presents the complete per-sector, per-regime economic validation
tables for Models A, B and C.
